%% =====================================================================
%%  T-02-04  Effortlessness
%%  -------------------------------------------------------------------
%%  One idea per file. Core is mandatory. Slots in canonical order.
%%  Max 700 words. No layout commands. No filler. New source material
%%  for this entry goes in sources/<BOOK>/T-02-04.tex -- never by
%%  rewriting this file (docs/RULES.md R0.1).
%% =====================================================================
%% @kind: topic
%% @id: T-02-04
%% @title: Effortlessness
%% @subtitle: Conceal the work or the work becomes the story
%% @chapter: c02
%% @order: 40
%% @status: stable
%% @sources: B3
%% @updated: 2026-09-19

\topic{T-02-04}{Effortlessness}{Conceal the work or the work becomes the story}

\begin{core}
An accomplishment that looks effortless raises your standing; the same accomplishment with the effort visible invites inspection, envy and imitation. What is concealed is not the result but the labour and the artifice behind it.
\end{core}
\begin{mechanism}
Visible effort invites two unwanted responses. It makes the result look attainable, which lowers the standing it buys you, and it exposes the method, which is the reusable part --- a method others can see is a method others can copy or counter. Effortlessness also disarms scrutiny: people audit what looks strained and accept what looks natural. The discipline is therefore presentational and informational at once.
\end{mechanism}
\begin{tells}
  \item You have described how hard something was, unprompted.
  \item You have explained a method to someone who has no need of it.
  \item Your preparation is visible in the delivery --- notes, hedging, over-structuring.
  \item You feel the need to be credited for the labour rather than the outcome.
\end{tells}
\begin{moves}
  \item Show the result, not the process.
  \item Rehearse until the delivery has no seams, then present it as though it were ordinary.
  \item Keep your methods to yourself, including from allies who have no use for them.
  \item When asked how you did it, answer at one level of abstraction above the truth.
\end{moves}
\begin{counters}
  \item Ask for the method, not the result. A person who cannot or will not describe how something was produced is protecting something.
  \item Reproduce the work yourself at small scale. Effortlessness does not survive an attempt.
  \item Watch what someone does not claim credit for; the gap between visible and invisible effort is where the artifice lives.
\end{counters}
\begin{cost}
It makes you un-teachable, which is expensive: a person who conceals method cannot delegate well, cannot build a team, and cannot be replaced --- which sounds like an advantage until the workload arrives. It also isolates, since the persona has to be maintained with everyone. And it fails outright where the audience is expert; effortlessness in front of people who know the craft reads as concealment and buys resentment instead of standing.
\end{cost}

\sources{B3}
\seesources
\seealso{T-09-01, T-13-04}
